% =====================================================================
%  前置部分：译校说明 + 作者前言（まえがき，原书 p.v–x）
% =====================================================================
\chapter*{译校说明}
\addcontentsline{toc}{chapter}{译校说明}
\markboth{译校说明}{译校说明}

\section*{一、底本}

本书译自深谷賢治（Kenji Fukaya）《シンプレクティック幾何学》，岩波書店，1999 年初版；2008 年由同一出版社发行单行本（本文内容与讲座版相同，仅订正排印错误，并追加“単行本化にあたっての追記”）。中译本以 2008 年单行本为底本，原书正文 372 页，另附前言、后记、追记、参考文献与索引。

\section*{二、体例}

\begin{enumerate}[label=\arabic*.,leftmargin=2.2em,itemsep=2pt]
  \item \textbf{页码对照}：正文中以 \texttt{\% [原文 p.N]} 标注各段所对应的原书页码；图注末尾标“（原书 p.N）”。原书页码 = PDF 页码 − 10。
  \item \textbf{公式与陈述编号}：原书公式按“章.序号”编号（如 (2.14)），各类陈述（定义、定理、引理、命题、推论、例、注、问题、猜想等）在每章内统一连续编号（如“引理 2.30”“定义 2.31”）。中译本一律照录原编号，不改动、不重排。
  \item \textbf{图}：原书插图均自扫描页裁出，图号与图注照原书；个别示意图若裁图失真，则以文字说明其内容。
  \item \textbf{译校批注}：以源码注释 \texttt{\% [Tutor Note: ...]} 形式给出，不进入正文排版，全部汇集于书末“附录：译校批注”，按章节排列并注明原书页码。批注只写三类内容：易错预警（符号、指标、约定）、作者省略论证的点拨、背景直觉。原书疑似笔误一律照原样排，只在批注中指出并给出正确形式。
  \item \textbf{本节重点与章末评注}：每节开头的“本节重点”方框与每章末尾的“本章回顾与评注”方框均为中译本所加，非原书内容。全书可用 \texttt{main.tex} 中的 \texttt{\textbackslash showtakeawayfalse} 与 \texttt{\textbackslash showchapterreviewfalse} 一键隐藏。
  \item \textbf{参考文献}：照录原书，不译条目。
\end{enumerate}

\section*{三、术语原则}

术语以“附录：术语对照表”为准。原则是：中文学界有稳定、无歧义通行译名者用中文；无稳定译名或硬译反碍对照文献者保留原文（如 operad、Fukaya 范畴、Floer 同调、Gromov–Witten 不变量等专名与缩写）。人名一律不译。日文与中文用法有别者按中文习惯处理，例如：日文“体”（代数学）译“域”，“補題”译“引理”，“系”译“推论”，“注意”译“注”，“概複素構造”译“近复结构”，“概正則曲線”译“伪全纯曲线”。

\section*{四、声明}

本书为学术研读用译本，不作商业用途。译文与批注中的错误由译者负责，恳请读者指正；原书的一切著作权归作者与岩波書店所有。

\clearpage

% =====================================================================
\chapter*{前言}
\addcontentsline{toc}{chapter}{前言（原书 p.v–x）}
\markboth{前言}{前言}

本书是辛几何的概说，以 Gromov 开创的、用伪全纯曲线研究整体辛几何的部分为中心。不过，我也尽量多触及除此以外的话题。

辛几何诞生于分析力学。20 世纪分析力学的发展有两条脉络：一是从三体问题等发展起来的动力系统研究，二是从量子力学走向量子场论的潮流\footnote{把这两条分开来看，当然是一面之见。}。本书的着重点可以说偏于后者，但也叙述了与前者相关的内容。

在辛流形整体性质的研究中，以辛流形分类等为中心的部分，近来常被称为辛拓扑。本书的着重点与辛拓扑有所偏离，本书并非辛拓扑的概说书。（本书着重点如何安排，背后有 [128]、[115]、[118] 等所述及的梦想。）

作为伪全纯曲线理论基础的、用泛函分析式偏微分方程论构造伪全纯曲线模空间并证明其基本性质——这些连同技术细节在内的完整叙述，本书并未给出。但要点我打算在第 2 章说明。关于伪全纯曲线模空间的构造，近来已有很多著作（[25]、[184]、[253]）。

辛几何近来涉及的重要话题之一是镜像对称。本书为此花了相当篇幅，但镜像对称的许多重要侧面本书并未触及。镜像对称同时与辛几何、代数几何·复几何相关。关于镜像对称的解说，大多站在理论物理学（弦理论）的立场，或站在代数几何学·复几何的立场（从代数几何立场出发的详细解说最近已经出版 [64]）。

本书当然是从辛几何的立场来解说的，对代数几何·复几何应当讨论的内容没有深入太多。特别是，对于镜像对称研究的一个中心——Calabi–Yau 流形上有理曲线的计数——本书没有深入。

本书所述及的镜像对称的侧面，与属于代数几何·复几何的侧面相比，有许多尚未成熟的部分；但我希望在这有限的篇幅内写出一本有特色的书，结果便是把与笔者正在研究的东西相近的部分（包括正在进行中的工作）作了某种程度的叙述。

辛几何从天体力学、动力系统、奇点论、流体力学等开始，直到微局部分析、可积系统、代数几何、粒子物理、李群表示论、非交换几何，与数学及其邻近诸科学的非常多的领域相关。本书动笔之后，“这个也得查”“那个也想学”的事情不断增加，笔者研究室的桌上总是堆着三四十本书和一百多篇论文，不停地新旧交替。其中一部分引用在本书中，但也有些只是装样子，其实是束之高阁而引用的。希望读者能按照各自的兴趣，继续学习与研究。

各章内容如下。第 1 章是导论，是笔者对辛几何在 20 世纪几何学整体中所处位置的一家之言。

第 2 章起是正文。第 2 章叙述辛流形的定义与辛几何最初的两个基本定理——Darboux 定理与 Moser 定理——之后，说明伪全纯曲线的方法。§2.4、§2.5、§2.6 阐述考察伪全纯曲线的重要性，以及伪全纯曲线模空间的构造与研究的基本想法。暂且作为目标的定理 2.36，是关于不可积的近复结构情形下 $\CP^2$ 上伪全纯曲线的存在定理。

与庞大的工具相比，这个定理乍看并不像大定理。首先要解说工具本身，这是目的所在；而在 §2.7 中，利用这个定理（的亲属），叙述宣告整体辛几何诞生的 Gromov 定理，即不可压缩定理。在不可压缩定理中，考虑不可积的近复结构是本质性的，这与复几何（凯勒流形的几何学）不同，强烈体现了辛几何的特征。

第 3 章是辛商（或者说辛约化）的概说。这方面也有许多著作（[24]、[167]、[169]、[200]）。本书叙述得比较详细的是动量映射的 Morse 理论。此外，作为例子，较多讨论了曲面基本群的表示空间。辛商与动量映射一直与李群表示论的研究融为一体，这方面的例子很多，但本书没有过多涉及。理由之一是笔者学识不足；另一是本书的着重点在于量子场论、特别是其与整体几何深刻相关的部分，而曲面基本群的表示空间从这个意义上也是重要的例子。基本群的表示空间，是全体联络所成空间的辛商。这类无穷维的例子还有很多，整体辛几何的视角今后也将是发挥威力的方向。

第 4 章讨论辛几何的重要概念——拉格朗日子流形。前两节的目的是尽量列举例子与构造法，展示拉格朗日子流形与多少数学相关。§4.3 解说几何量子化。几何量子化同样是与表示论深刻相关而发展起来的领域，实例也多与表示论有关；但出于与前面相同的理由，本书选择了与曲面基本群表示空间相关的例子。§4.4 解说 Maslov 指数。Maslov 指数也是与多方面相关的话题。它作为示性类自然地与拉格朗日配边理论相关，对此本书叙述得相当详细。此外，对于与最近发展相关的重要侧面——与伪全纯圆盘模空间的关系，也作了叙述。另一个重要侧面——与微局部分析的关系，遗憾的是限于篇幅割爱了。请参阅 [34]、[165] 等。

第 5 章概述辛几何中的 Floer 同调理论。辛几何中出现的 Floer 同调（暂且）有两种：前三节叙述周期哈密顿系统的情形，§5.4 叙述拉格朗日子流形的情形（最易处理的情形）。遗憾的是未能把必要的分析细节讲到底，但我想已展示了 Floer 同调理论的基本问题与想法。

第 6 章处理镜像对称。系统地用伪全纯曲线组装起来的，是 Gromov–Witten 不变量的理论。最初三节是其概要。仅基本构造全部写出就会变得冗长，因此省略了许多分析细节。（分析的基本想法在第 2 章。）此外，关于 Gromov–Witten 不变量的计算方法，本书讲得不多。前面引用过的、关于镜像对称代数几何侧面的概说 [64] 对计算方法有详细论述。后三节叙述形变理论（模空间理论）及其与辛几何的关系。镜像对称与模空间发生关系，首先是通过周期映射，但本书没有说明周期映射。镜像对称与周期映射的关系在 [64] 中也很详细。（计算法与周期映射都属于代数几何的话题，故避开；同时也尽量避免了与 [64] 重复。）作为替代，这三节叙述基于同调（同伦）代数的抽象形变理论，并给出其在镜像对称与辛几何中的应用。§6.4 叙述同调代数的部分（连同来自复流形的动机），§6.5 叙述其应用——Kontsevich 等人关于形变存在性的两个定理，即 Calabi–Yau 流形的扩展形变空间的构造\footnote{实际上这与周期映射的理论关系很深。}，以及相应于泊松结构的形变量子化的存在性。最后一节叙述大部分仍作为猜想残留的同调镜像对称。主要目的是导出数学上严格形式的猜想，并给出阿贝尔流形情形下猜想成立的具体例子。这一节的一半，是预定与本书同时期出版的、Kontsevich、Oh、太田、小野与笔者共著的 [122] 的概要。

本书各章有相当部分可以互相独立地阅读。阅读第 3 章与第 4 章（§4.4 最后部分除外）时，第 2 章 §2.3 以后不需要。又，第 4 章的大部分与第 2 章独立；用到第 3 章结果的地方只有 §4.3 的一部分；§4.4 与 §4.3 独立。第 5 章的 §5.1、§5.2、§5.3 与第 4 章独立，第 3 章中只用 §3.1、§3.2。§5.4 则用到第 4 章的 §4.1 与 §4.4。第 5 章与第 6 章也几乎独立；只要看一下普遍 Novikov 环的定义 5.44，就可以在第 2 章之后直接进入第 6 章。不过 §6.6 是例外，与第 3、4、5 章的大部分相关。

本书各处（特别是第 5、6 章）以“注”的形式，写明了假定读者具备相当预备知识。初读的读者不能理解这些的大部分，是理所当然的。遇到不懂的词，尽可快快跳过。这些注，是笔者希望读者各自对其中的一两处发生兴趣、并以所引文献为线索推进学习与研究而写的。

笔者开始学习辛几何是在六七年前，起因是在研究生讨论班上提出“来做辛几何吧”。此后又以硕士课程、博士课程讨论班的名义，请京都大学与东京大学的研究生们读各种论文并讲解。这些构成了本书的主要部分。此外，在研究集会系列 “Surveys in Geometry”（特别是 [129]、[130]、[131] 那几回）及其准备过程中，也向许多人学到了各种东西。

如此，从正在学习的学生到辛几何的代表性研究者，若没有许多人的教导，本书是写不出来的。以下按五十音顺序记下一部分人的名字，以表谢意。

赤穂まなぶ、板垣裕也、江口彻、Y.G.Oh、太田启史、股岛靖文、小野薫、金子真隆、龟谷幸夫、M. Gromov、牛肠彻、今野宏、M. Kontsevich、高仓树、中岛启、西纳武男、桥本义武、林正人、古田干雄……\footnote{以上人名照原书五十音顺序列出；汉字写法依原书，中国读者可参看各人论文。}

最后，对在笔者大幅拖延执笔、超出页数等给您们添了麻烦的情况下，仍然帮助本书写到最后的岩波书店编辑部各位，谨表谢意。

1999 年 5 月\quad 深谷賢治

\medskip
\noindent\textbf{追记}

本书是把原为岩波讲座《现代数学的展开》之一分册的《辛几何》单行本化的产物。单行本化时，正文内容一仍讲座之旧，仅订正排印错误。不过，在后记之后加入了“单行本化之际的追记”，对正文作了补充与订正。

2008 年 10 月
